\section{Sampling and orthogonal packing}\label{r:sec:packing}

Continue with the inclusion \(A\subseteq D\) from
Section~\ref{r:sec:weighted}. The matrix compressions from
Proposition~\ref{r:weighted:family} allow
us to sample a projection without making the sampled corner central.
We isolate the precise estimate before carrying out the packing.

\begin{lemma}[Sampling a matrix block]\label{r:packing:sampling}
Let \(q\in D\) be a nonzero projection and let \(U\in M_{1,k}(D)\)
satisfy \(U^*U=1_k\otimes q\). Put \(p=UU^*\).
Suppose \(x\in\U(D)\), \(0<\eta<1\), and
\[
 \dist(U^*xU,M_k(\C)\otimes q)<\eta.
\]
For a projection \(e\leq q\), set \(y=U(1_k\otimes e)U^*\).
Then, for \(0\leq c\leq1\), every \(\tau\in\T(D)\), and
\(\delta_\tau=|\tau(y)-c\tau(p)|\),
\begin{align}
 |\tau(yx)-c\tau(px)|
 &\leq\eta\big(\tau(y)+c\tau(p)\big)+2\delta_\tau,
 \label{r:eq:sampling-linear}\\
 \left|\normtwo{[x,y]}{\tau}^2
        -c\normtwo{[x,p]}{\tau}^2\right|
 &\leq6\eta\big(\tau(y)+c\tau(p)\big)+10\delta_\tau.
 \label{r:eq:sampling-energy}
\end{align}
\end{lemma}

\begin{proof}
Write \(T=U^*xU\), \(P_e=1_k\otimes e\), and choose \(T_0\in M_k(\C)\)
with \(\norm{T-T_0\otimes q}<\eta\).
Since \(T\) is a contraction, \(\norm{T_0}\leq2\).
The map \(a\mapsto UaU^*\) identifies the relevant matrix corner with
its image in \(D\), and therefore
\[
 \tau(yx)=(\Tr\otimes\tau)(P_eTP_e),\qquad
 \tau(|yxy|^2)=(\Tr\otimes\tau)(|P_eTP_e|^2).
\]
The scalar matrix \(T_0\otimes q\) commutes with \(P_e\), whether or not
\(e\) commutes with any entry of \(T\).
Thus replacing \(P_eTP_e\) by \(T_0\otimes e\) changes these two traces
by at most \(\eta\tau(y)\) and \(3\eta\tau(y)\), respectively.
For the latter estimate use
\[
 \norm{R^*R-S^*S}\leq(\norm R+\norm S)\norm{R-S},
 \quad R=P_eTP_e,\quad S=T_0\otimes e,
\]
and the fact that the trace of the corner unit is
\((\Tr\otimes\tau)(P_e)=k\tau(e)=\tau(y)\).
The same estimates hold with \(e\) replaced by \(q\).

Set
\[
 \alpha=\frac{\Tr(T_0)}{k},\qquad
 \beta_0=\frac{\Tr(T_0^*T_0)}{k}.
\]
Then \(|\alpha|\leq2\) and \(0\leq\beta_0\leq4\).
The scalar reference traces for \(y\) are
\(\alpha\tau(y)\) and \(\beta_0\tau(y)\); for \(p\), they are
\(\alpha\tau(p)\) and \(\beta_0\tau(p)\).
This proves \eqref{r:eq:sampling-linear}.
For the second estimate, the exact identity
\begin{equation}\label{r:eq:unitary-energy}
 \normtwo{[x,f]}{\tau}^2
 =2\tau(f)-2\tau(|fxf|^2)
 \qquad(f\text{ a projection},\ x\text{ unitary})
\end{equation}
follows by expansion and traciality. Applying it to \(y,p\) and using
the squared compression estimates proves \eqref{r:eq:sampling-energy};
the mass-discrepancy coefficient is bounded by
\(2+2\beta_0\leq10\).
In particular, \eqref{r:eq:unitary-energy} also accounts for the part of
\(x\) that leaves the larger block \(p\).
\end{proof}

\begin{lemma}[Energy of orthogonal sums]\label{r:packing:orthogonal}
If \(t_1,\ldots,t_M\) are mutually orthogonal projections and
\(z=\sum_a t_a\), then for \(x\in D\) and \(\tau\in\T(D)\),
\[
 \normtwo{[x,z]}{\tau}^2
 \leq\sum_a\normtwo{[x,t_a]}{\tau}^2.
\]
The same bound holds for a sum over any subfamily.
\end{lemma}

\begin{proof}
The two off-diagonal corners of the commutator are orthogonal in the
tracial 2-seminorm, so its square is
\[
 \normtwo{(1-z)xz}{\tau}^2+\normtwo{zx(1-z)}{\tau}^2.
\]
Splitting the \(z\)-corner into the orthogonal \(t_a\)'s writes each
term as a sum of squared norms. Since \(1-z\leq1-t_a\), the respective
summands are bounded by
\(\normtwo{(1-t_a)xt_a}{\tau}^2\) and
\(\normtwo{t_ax(1-t_a)}{\tau}^2\).
\end{proof}

\begin{proof}[Proof of Theorem~\ref{r:thm:finite}]
Fix a finite \(K\subseteq\U(A)\).
We construct orthogonal projections \(P,S\) such that \(P\) nearly commutes
with \(K\), their coefficient traces nearly agree, and \(P+S\) nearly
fills the unit in trace. The identity
\[
 2\tau(Px)-\tau(x)
 =\tau((P-S)x)-\tau((1-P-S)x)
 \qquad(x\in K)
\]
then gives the required halving estimate. We obtain \(P,S\) as two
labelled sums, sampling equal proportions for the two signs in each
round and moving the samples into the unoccupied corner.
We use \(H>0\) and \(N\in\N\), with \(\eta=H/N<1\).
First \(H\) is chosen to reduce the unoccupied mass; then \(N\) is
increased at fixed \(H\) to reduce the errors.

\paragraph{Preparing the averaging maps.}
When an averaging map acts on the remaining corner, the projections
inserted earlier must already approximately commute with its coefficient
contractions. Since the map is chosen independently of its input
projection, we can arrange these coefficient tests in advance, working
backwards through the rounds.
Construct finite sets of unitaries
\[
 K_1\supseteq K_2\supseteq\cdots\supseteq K_N=K
\]
backward. Having chosen \(K_i\), choose the map \(\Phi_i\) from
Proposition~\ref{r:weighted:family} for \(K_i,\eta\).
If \(i>1\), express each of its coefficient contractions \(v_l\) as a
linear combination of unitaries with coefficient absolute sum at most
two, and include these unitaries, together with \(K_i\), in \(K_{i-1}\).

\paragraph{The induction.}
We insert finitely many projections in each of \(N\) rounds.
All inserted projections will be mutually orthogonal and each will
carry a label \(+\) or \(-\).
Let \(Z_i\) be their cumulative sum after round \(i\), with \(Z_0=0\).
Set
\begin{equation}\label{r:eq:energy-recurrence}
 E_0=0,\qquad E_i-E_{i-1}=1000\eta(E_{i-1}+\eta).
\end{equation}
We maintain the bound
\begin{equation}\label{r:eq:energy-induction}
 \sum_{t\text{ inserted through round }i}
       \normtwo{[x,t]}{\tau}^2\leq E_i
 \quad(x\in K_i,\ \tau\in\T(D)).
\end{equation}
In particular, Lemma~\ref{r:packing:orthogonal} gives the same bound for
the squared commutator of any subfamily sum.

In round \(i\), let \(r=1-Z_{i-1}\) and choose the weighted family from
Proposition~\ref{r:weighted:family} for \(\Phi_i,r\).
For every \(\tau\in\T(D)\),
\begin{equation}\label{r:eq:old-overlap}
 \tau(Z_{i-1}\Phi_i(r))
 =\sum_l\lambda_l\normtwo{Z_{i-1}v_lr}{\tau}^2
 \leq4E_{i-1}.
\end{equation}
For \(i=1\) the left-hand side is zero. Otherwise each unitary in the
chosen decomposition of \(v_l\) lies in \(K_{i-1}\), so the induction
and Lemma~\ref{r:packing:orthogonal} give
\(\normtwo{[Z_{i-1},v_l]}{\tau}\leq2\sqrt{E_{i-1}}\).
Also \(Z_{i-1}v_lr=[Z_{i-1},v_l]r\), proving
\eqref{r:eq:old-overlap}.

\paragraph{One round of sampling and insertion.}
Within this round we suppress \(i\) from the block and signed-index
notation, and process the finite signed index set \(a=(j,\pm)\) in any
fixed order.
Both indices for \(j\) have the target proportion
\[
 c_a=\eta w_j/2.
\]
Let \(z_a\) be the occupied projection just before index \(a\), starting
with \(Z_{i-1}\). We choose
\[
 y_a=U_j(1_{k_j}\otimes e_a)U_j^*,\qquad e_a\leq q_j,
\]
so that, uniformly over all \(\tau\in\T(D)\),
\begin{align}
 \sum_a|\tau(y_a)-c_a\tau(p_j)|&\leq\eta^2,
       \label{r:eq:sample-mass-errors}\\
 \sum_a|\tau(z_ay_a)-c_a\tau(z_ap_j)|&\leq\eta^2.
       \label{r:eq:sample-overlap-errors}
\end{align}
Here \(j\) always denotes the block index of \(a\).
To make these choices precise, assign positive budgets
\(\eps_a\) with \(\sum_a\eps_a\leq\eta^2\) before processing the family.
At step \(a\), apply Corollary~\ref{r:proj:tests} with proportion \(c_a\),
tests \(q_j\) and
\[
 h_a=\sum_{d=1}^{k_j}(U_j^*z_aU_j)_{dd}\in q_jDq_j,
\]
and tolerance \(\eps_a/k_j\).
Traciality gives
\[
 \tau(y_a)=k_j\tau(e_a),\qquad
 \tau(z_ay_a)=\tau(e_ah_a),\qquad
 \tau(z_ap_j)=\tau(q_jh_a).
\]
Thus both errors at this step are at most \(\eps_a\).
The test \(h_a\), which satisfies \(0\leq h_a\leq k_jq_j\), is known
when it is used. This application may refine \(q_j\) internally, but
only returns a subprojection \(e_a\); it does not change the fixed
weighted family.

Use Lemma~\ref{r:proj:move} to replace \(y_a\) by a projection
\(t_a\leq1-z_a\), retaining its sign, such that
\begin{equation}\label{r:eq:movement}
 \tau(y_a)-3\tau(z_ay_a)\leq\tau(t_a)\leq\tau(y_a),
 \qquad
 \normtwo{t_a-y_a}{\tau}^2\leq5\tau(z_ay_a).
\end{equation}
Update the occupied projection to \(z_a+t_a\).
At the end of the round this is \(Z_i\).

\paragraph{Mass and overlap.}
For the estimates in this round, fix \(\tau\in\T(D)\), and put
\[
 C_i(\tau)=\sum_a c_a\tau(p_j),\quad
 Y_i(\tau)=\sum_a\tau(y_a),\quad
 L_i(\tau)=\sum_a\tau(z_ay_a).
\]
These are the target mass, sampled mass, and overlap with the occupied
projection in this round. All subsequent bounds hold for every \(\tau\)
with the same constants; we suppress \(\tau\) in these three quantities.
Equations \eqref{r:eq:diagonal-mass} and \eqref{r:eq:block-mass} give
\[
 \eta\tau(r)-2\eta^3\leq C_i\leq\eta(1+\eta^2)\leq2\eta.
\]
Using \eqref{r:eq:sample-mass-errors}, we obtain
\begin{equation}\label{r:eq:round-mass}
 \eta\tau(r)-3\eta^2\leq Y_i\leq3\eta.
\end{equation}
For the overlap, first use \(z_a\leq Z_i\), then apply the single
weighted-operator estimate \eqref{r:eq:block-mass}:
\begin{align}
 L_i
 &\leq\eta^2+\eta\sum_jw_j\tau(Z_ip_j)\notag\\
 &\leq\eta^2+\eta\big(\eta^2+\tau(Z_i\Phi_i(r))\big)\notag\\
 &\leq\eta^2+\eta\big(\eta^2+4E_{i-1}+Y_i\big)\notag\\
 &\leq4\eta E_{i-1}+5\eta^2.
 \label{r:eq:round-overlap}
\end{align}
In the third line we used \eqref{r:eq:old-overlap},
\(0\leq\Phi_i(r)\leq1\), and
\(\tau(Z_i-Z_{i-1})\leq Y_i\), the last following from
\eqref{r:eq:movement}.
In particular,
\begin{equation}\label{r:eq:overlap-summable}
 L_i\leq E_i-E_{i-1}.
\end{equation}

\paragraph{Commutators and linear tests.}
For \(x\in K_i\), sum Lemma~\ref{r:packing:sampling} over this round.
Equations \eqref{r:eq:sample-mass-errors} and \eqref{r:eq:round-mass}
give
\begin{align}
 \sum_a|\tau(y_ax)-c_a\tau(p_jx)|
 &\leq\eta(Y_i+C_i)+2\eta^2
 \leq7\eta^2,\label{r:eq:round-linear}\\
 \sum_a\normtwo{[x,y_a]}{\tau}^2
 &\leq\eta\sum_jw_j\normtwo{[x,p_j]}{\tau}^2
       +6\eta(Y_i+C_i)+10\eta^2\notag\\
 &\leq46\eta^2\leq60\eta^2.
 \label{r:eq:round-sample-energy}
\end{align}
Here the first term in the second estimate is at most \(6\eta^2\)
by \eqref{r:eq:block-energy}. Passing to the inserted projections,
\eqref{r:eq:movement} implies
\begin{align}
 \sum_a\normtwo{[x,t_a]}{\tau}^2
 &\leq2\sum_a\normtwo{[x,y_a]}{\tau}^2
        +8\sum_a\normtwo{t_a-y_a}{\tau}^2\notag\\
 &\leq120\eta^2+40L_i\notag\\
 &\leq160\eta E_{i-1}+320\eta^2
 \leq E_i-E_{i-1}.
 \label{r:eq:round-inserted-energy}
\end{align}
Together with \(K_i\subseteq K_{i-1}\), this proves
\eqref{r:eq:energy-induction}.
The mass left after the round satisfies
\begin{equation}\label{r:eq:remaining-mass}
 \tau(1-Z_i)\leq(1-\eta)\tau(1-Z_{i-1})+3\eta^2+3L_i,
\end{equation}
by \eqref{r:eq:round-mass} and the first inequality in
\eqref{r:eq:movement}.
All sums in this argument are controlled by weighted projection masses
or energies. No bound on the raw sum \(\sum_jw_j\), or on the number of
blocks, has been used.

\paragraph{Completing the packing.}
Let \(P,S\) be the sums of all inserted projections labelled \(+\) and
\(-\), respectively. Then \(P+S=Z_N\), and for
\(x\in K\), \(\tau\in\T(D)\),
\begin{align}
 \normtwo{[x,P]}{\tau}^2&\leq E_N,\label{r:eq:final-energy}\\
 \tau(1-Z_N)&\leq e^{-H}+3N\eta^2+3E_N,
       \label{r:eq:final-remainder}\\
 |\tau(Px)-\tau(Sx)|
 &\leq7N\eta^2+2\sqrt{15HE_N}.
       \label{r:eq:final-balance}
\end{align}
The first follows from Lemma~\ref{r:packing:orthogonal}.
For the second, iterate \eqref{r:eq:remaining-mass}, use
\((1-\eta)^N\leq e^{-H}\), and sum
\eqref{r:eq:overlap-summable}.
For the third, the \(+\) and \(-\) targets \(c_a\tau(p_jx)\) cancel
for every block in every round. Their total sampling error is at most
\(7N\eta^2\), by \eqref{r:eq:round-linear}.
The remaining displacement error is, by Lemma~\ref{r:proj:move} and
Cauchy--Schwarz over all signed indices,
\begin{align*}
 \sum_{i,a}|\tau((t_a-y_a)x)|
 &\leq2\sum_{i,a}\sqrt{\tau(y_a)}
                          \normtwo{t_a-y_a}{\tau}\\
 &\leq2\sqrt{\sum_iY_i}\sqrt{5\sum_iL_i}
 \leq2\sqrt{15HE_N}.
\end{align*}

Solving \eqref{r:eq:energy-recurrence} gives
\[
 E_N=\eta\big((1+1000\eta)^N-1\big)\leq\eta e^{1000H},
 \qquad N\eta^2=H\eta.
\]
Since \(x\) is a contraction,
\(|\tau((1-Z_N)x)|\leq\tau(1-Z_N)\).
Equations \eqref{r:eq:final-remainder} and \eqref{r:eq:final-balance}
therefore yield
\begin{equation}\label{r:eq:halving-error}
 \left|\tau(Px)-\tfrac12\tau(x)\right|
 \leq\tfrac12e^{-H}+5H\eta+\tfrac32E_N+\sqrt{15HE_N}.
\end{equation}
Choose \(H\) large enough to make the first term small. With this \(H\)
fixed, let \(N\) increase through integers and put \(\eta=H/N\).
Then every other term in \eqref{r:eq:halving-error}, and
\(\sqrt{E_N}\) in \eqref{r:eq:final-energy}, tends to zero.
This proves the two desired uniform estimates.
\end{proof}
